\section{Mean Bias as the Dominant Source of Activation Outliers}
\label{sec:activation_outliers}

We now address its sharp structural consequence: the production of extreme elementwise activations that destabilize low-bit quantization and adaptive optimization. Because the dynamic range of low-bit quantization is strictly governed by entrywise maxima (the $L_\infty$ norm), we perform an elementwise mechanical decomposition to precisely locate the geometric source of this instability.

To attribute outlier energy, we decompose the activation matrix $X \in \mathbb{R}^{N \times H}$ into three mutually orthogonal components. Let $\mu = \frac{1}{N}X^\top \mathbf{1} \in \mathbb{R}^H$ be the feature-wise mean vector, and let $\tilde{X} = X - \mathbf{1}\mu^\top$ be the column-centered activation matrix. We decompose $X$ as:
\begin{equation}
X = \underbrace{\mathbf{1}\mu^\top}_{M (\text{mean})} + \underbrace{U_k \Sigma_k V_k^\top}_{X_{\text{spike}}} + \underbrace{X_{\text{tail}}}_{\text{residual}},
\end{equation}
where $X_{\text{spike}}$ is the truncated rank-$k$ singular value decomposition of the centered matrix $\tilde{X}$ onto its top $k = \lfloor 0.01 H \rfloor$ singular directions, and $X_{\text{tail}} = \tilde{X} - X_{\text{spike}}$ is the remaining orthogonal residual. Because $X_{\text{spike}}$ and $X_{\text{tail}}$ are linear projections of the centered matrix, they maintain zero column means (i.e., they are orthogonal to the all-ones vector $\mathbf{1}$). Therefore, the decomposition is strictly orthogonal in the Frobenius inner product, yielding exact energy separation:
\begin{equation}
\|X\|_F^2 = \|M\|_F^2 + \|X_{\text{spike}}\|_F^2 + \|X_{\text{tail}}\|_F^2.
\end{equation}

\paragraph{Outlier Attribution.} 
Define the outlier set $\mathcal{E}_{\text{top}}$ as the indices $(i,j)$ of the top $0.1\%$ of entries ranked by absolute magnitude $|X_{ij}|$. To determine which geometric component drives these extreme values, we define the squared-magnitude share for each $(i,j) \in \mathcal{E}_{\text{top}}$:
\begin{equation}
\rho_{ij}^{(\text{mean})} = \frac{M_{ij}^2}{X_{ij}^2}, \quad \rho_{ij}^{(\text{spike})} = \frac{(X_{\text{spike}})_{ij}^2}{X_{ij}^2}, \quad \rho_{ij}^{(\text{tail})} = \frac{(X_{\text{tail}})_{ij}^2}{X_{ij}^2}.
\end{equation}

\textit{Remark:} While the decomposition is strictly orthogonal over the full matrix, individual elementwise squares technically contain minor cross-terms (e.g., $2M_{ij}(X_{\text{spike}})_{ij}$). Nevertheless, measuring these primary component ratios provides a highly accurate and stable proxy for identifying the dominant driver of extreme values.

\begin{table}[htbp]
\centering
\caption{Outlier contribution shares (top 0.1\% entries, squared magnitude).}
\label{tab:outlier_shares}
\begin{tabular}{lccc}
\toprule
Layer / Step & Mean (avg) & Spike (avg) & Tail (avg) \\
\midrule
Layer 3 (10k)   & 0.60  & 0.07  & 0.01  \\
Layer 3 (180k)  & 0.98  & 0.004 & 0.004 \\
Layer 21 (10k)  & 0.004 & 0.038 & 0.70  \\
Layer 21 (180k) & 0.70  & 0.038 & 0.015 \\
\bottomrule
\end{tabular}
\end{table}

\begin{figure}[htbp]
\centering
\begin{subfigure}{0.24\textwidth}
    \centering
    \includegraphics[width=\linewidth]{./figures/layer3_step010000_activation.pt_outlier_hist.png} % Replace with your actual image file
    \caption{Layer 3, 10k}
\end{subfigure}\hfill
\begin{subfigure}{0.24\textwidth}
    \centering
    \includegraphics[width=\linewidth]{./figures/layer3_step010000_activation.pt_outlier_hist.png} % Replace with your actual image file
    \caption{Layer 3, 180k}
\end{subfigure}\hfill
\begin{subfigure}{0.24\textwidth}
    \centering
    \includegraphics[width=\linewidth]{./figures/layer21_step010000_activation.pt_outlier_hist.png} % Replace with your actual image file
    \caption{Layer 21, 10k}
\end{subfigure}\hfill
\begin{subfigure}{0.24\textwidth}
    \centering
    \includegraphics[width=\linewidth]{./figures/layer21_step180000_activation.pt_outlier_hist.png} % Replace with your actual image file
    \caption{Layer 21, 180k}
\end{subfigure}
\caption{Activation outlier contribution shares (top 0.1\% entries). Mean dominance progressively emerges across layers and training.}
\label{fig:outlier_shares}
\end{figure}

Table~\ref{tab:outlier_shares} and Figure~\ref{fig:outlier_shares} summarize these statistics across layers and training stages. Two distinct structural transitions emerge. In shallow layers early in training, mean bias already accounts for the vast majority of extreme entries. As training proceeds, shallow-layer outliers become almost entirely mean-dominated. In deeper layers, the spike variance initially dominates outliers, but by late training, mean bias transitions into the principal driver across the entire network.

\paragraph{High-Dimensional Extreme-Value Amplification.} 
We now mathematically formalize why a rank-one mean shift inevitably dominates extreme elementwise values in high-dimensional settings, overpowering the variance.

\begin{theorem}[Elementwise Extreme Dominance]
\label{thm:extreme_dominance}
Let an activation coordinate be modeled as $X_{ij} = \mu_j + Z_{ij}$, where $\mu_j$ is a deterministic mean shift and $Z_{ij}$ is independent zero-mean sub-Gaussian noise with variance proxy $\sigma^2$. For a high quantization threshold $t > 0$ such that $t < |\mu_j|$, the probability of the coordinate magnitude exceeding the threshold satisfies:
\begin{equation}
\mathbb{P}(|X_{ij}| > t) \ge 1 - 2\exp\left(-\frac{(|\mu_j| - t)^2}{2\sigma^2}\right).
\end{equation}
\end{theorem}

\begin{proof}
By the reverse triangle inequality, $|X_{ij}| = |\mu_j + Z_{ij}| \ge |\mu_j| - |Z_{ij}|$. Therefore, if the ambient noise fluctuation satisfies $|Z_{ij}| < |\mu_j| - t$, it strictly guarantees that $|X_{ij}| > t$. 

Because $Z_{ij}$ is $\sigma^2$-sub-Gaussian, the probability of the complementary tail event is bounded by $\mathbb{P}(|Z_{ij}| \ge |\mu_j| - t) \le 2\exp\left(-\frac{(|\mu_j| - t)^2}{2\sigma^2}\right)$. Subtracting this from $1$ yields the lower bound. As the noise variance $\sigma^2$ becomes small relative to the margin $(|\mu_j| - t)^2$, this probability rapidly approaches $1$.

In stark contrast, if an activation coordinate is purely variance-driven ($X'_{ij} = Z_{ij}$ without a mean shift), exceeding the threshold $t$ requires a rare tail event, which is strictly bounded by standard sub-Gaussian concentration:
\begin{equation}
\mathbb{P}(|Z_{ij}| > t) \le 2\exp\left(-\frac{t^2}{2\sigma^2}\right).
\end{equation}
\end{proof}

This theorem formalizes a critical statistical divergence: a deterministic mean shift generates extreme values with constant probability $\mathcal{O}(1)$ (with respect to $H$), while purely variance-driven extremes occur with exponentially vanishing probability $\mathcal{O}(\exp(-t^2/\sigma^2))$. 

Because the dynamic range for blockwise low-bit tensor quantization is determined by the $L_\infty$ norm (the maximum scalar element in a block), the deterministic mean bias structurally stretches the quantization grid. This forces the grid to accommodate the rank-one drift, fatally starving the variance-driven semantic tail of bit-depth. Mitigating this mean bias therefore directly targets the geometric root of extreme-value instability, restoring resolution to the semantic signal.